% Options for packages loaded elsewhere
\PassOptionsToPackage{unicode}{hyperref}
\PassOptionsToPackage{hyphens}{url}
\PassOptionsToPackage{dvipsnames,svgnames,x11names}{xcolor}
%
\documentclass[
  11pt,
]{article}
\usepackage{amsmath,amssymb}
\usepackage{iftex}
\ifPDFTeX
  \usepackage[T1]{fontenc}
  \usepackage[utf8]{inputenc}
  \usepackage{textcomp} % provide euro and other symbols
\else % if luatex or xetex
  \usepackage{unicode-math} % this also loads fontspec
  \defaultfontfeatures{Scale=MatchLowercase}
  \defaultfontfeatures[\rmfamily]{Ligatures=TeX,Scale=1}
\fi
\usepackage{lmodern}
\ifPDFTeX\else
  % xetex/luatex font selection
\fi
% Use upquote if available, for straight quotes in verbatim environments
\IfFileExists{upquote.sty}{\usepackage{upquote}}{}
\IfFileExists{microtype.sty}{% use microtype if available
  \usepackage[]{microtype}
  \UseMicrotypeSet[protrusion]{basicmath} % disable protrusion for tt fonts
}{}
\makeatletter
\@ifundefined{KOMAClassName}{% if non-KOMA class
  \IfFileExists{parskip.sty}{%
    \usepackage{parskip}
  }{% else
    \setlength{\parindent}{0pt}
    \setlength{\parskip}{6pt plus 2pt minus 1pt}}
}{% if KOMA class
  \KOMAoptions{parskip=half}}
\makeatother
\usepackage{xcolor}
\usepackage[margin=2.2cm]{geometry}
\setlength{\emergencystretch}{3em} % prevent overfull lines
\providecommand{\tightlist}{%
  \setlength{\itemsep}{0pt}\setlength{\parskip}{0pt}}
\setcounter{secnumdepth}{-\maxdimen} % remove section numbering
\ifLuaTeX
  \usepackage{selnolig}  % disable illegal ligatures
\fi
\usepackage{bookmark}
\IfFileExists{xurl.sty}{\usepackage{xurl}}{} % add URL line breaks if available
\urlstyle{same}
\hypersetup{
  pdftitle={Draft: a(3,n) = n+2 (working notes)},
  colorlinks=true,
  linkcolor={Maroon},
  filecolor={Maroon},
  citecolor={Blue},
  urlcolor={Blue},
  pdfcreator={LaTeX via pandoc}}

\title{Draft: a(3,n) = n+2 (working notes)}
\author{}
\date{}

\begin{document}
\maketitle

\section{The Maximum Right-Isosceles-Free Subset of a Three-Row
Board}\label{the-maximum-right-isosceles-free-subset-of-a-three-row-board}

\subsection{Abstract}\label{abstract}

Let a(3,n) denote the largest number of lattice points that can be
chosen from a board with 3 rows and n columns so that no three chosen
points form a non-degenerate right isosceles triangle. We prove that

\textbf{a(3,n) = n + 2 for every n ≥ 4.}

The lower bound is an explicit construction valid for all n.~The upper
bound is proved by a potential-function argument on a finite automaton
of 4-column states. The finite part of the upper bound is checked by
computer, and the certificate is a short finite table that can be
verified independently. The argument is therefore computer-assisted,
though every reduction step is elementary.

\subsection{1. Setting}\label{setting}

A \emph{board with 3 rows and n columns} is the set of lattice points
(x, y) with 0 ≤ x ≤ n−1 and y ∈ \{0,1,2\}. A set S of points is
\emph{admissible} if it contains no three points forming a
non-degenerate right isosceles triangle, meaning a triple \{A, B, C\}
with a right angle at some vertex V and equal legs VA = VB ≠ 0.

Column j of a set S is encoded as the 3-bit mask c\_j = Σ\_\{y ∈ rows of
S in column j\} 2\^{}y, with values in \{0,\ldots,7\}. A sequence (c\_0,
\ldots, c\_\{n−1\}) of masks determines S. We write pop(c) for the
number of set bits, so the number of points in column j is pop(c\_j).

\subsection{2. Two elementary lemmas}\label{two-elementary-lemmas}

\textbf{Lemma 1 (horizontal extent).} Any right isosceles triangle with
vertices on the 3-row board has all its vertices within a horizontal
range of at most 4 columns. More precisely, its x-coordinates span at
most 4 consecutive integers.

\emph{Proof.} Let the right angle be at P, and let the legs be u = (p,
q) and w = ε(−q, p) with ε ∈ \{+1, −1\}, since w is a rotation of u by
±90°. The vertices, relative to P, have y-coordinates 0, q and εp. These
lie within an interval of length at most 2 (the board has 3 rows), and
that interval contains 0. Hence \textbar q\textbar{} ≤ 2 and
\textbar p\textbar{} ≤ 2. The x-coordinates relative to P are 0, p and
−εq, all lying in {[}−2, 2{]}. So the x-span is at most 4. ∎

\textbf{Lemma 2 (locality).} Let n ≥ 5. A set S on the 3×n board is
admissible if and only if every set of five consecutive columns of S is
admissible.

\emph{Proof.} Forward direction: restriction of an admissible set to a
sub-board is admissible. Backward direction: any triangle in S has
x-span at most 4 by Lemma 1, so it lies inside some five consecutive
columns. ∎

For n ≥ 5, Lemma 2 makes admissibility a \emph{local} condition on
windows of five columns, so the problem is a path problem in a finite
graph. The graph is described next.

\subsection{3. The state graph}\label{the-state-graph}

A \textbf{state} is an admissible 4-tuple of masks s = (c\_1, c\_2,
c\_3, c\_4), meaning that the 4-column board they encode is admissible.
A \textbf{transition} s → t, where t = (c\_2, c\_3, c\_4, c\_5), is
allowed when the 5-tuple (c\_1, \ldots, c\_5) is admissible. The
transition has \textbf{weight} pop(c\_5), the number of points added in
the new column. By Lemma 2, an admissible 3×n set with n ≥ 4 corresponds
exactly to a path of n−3 states, where consecutive states differ by an
allowed transition. The total number of points is

T = pop(c\_1) + pop(c\_2) + pop(c\_3) + pop(c\_4) + Σ over transitions
of their weights.

Computationally there are 443 states and 1736 admissible 5-tuples.

\subsection{4. Upper bound: a(3,n) ≤ n + 2}\label{upper-bound-a3n-n-2}

\textbf{Proposition 3 (certificate).} There is a function Φ from the set
of states to the integers \{0, 1, 2\} with the following properties.

\begin{enumerate}
\def\labelenumi{(\alph{enumi})}
\tightlist
\item
  Φ(s) ≥ 0 for every state s.
\item
  For every allowed transition s → t with new column c\_5: Φ(s) ≥
  (pop(c\_5) − 1) + Φ(t).
\item
  For every state s: pop(c\_1) + pop(c\_2) + pop(c\_3) + pop(c\_4) +
  Φ(s) ≤ 6.
\end{enumerate}

A concrete Φ is given in \texttt{phi3.json} (keys are states, values are
Φ). It is computed as the longest-path potential: Φ(s) is the supremum
over allowed paths starting at s of the sum of (weight − 1). This
supremum is finite, which is the computational content of the claim that
no cycle of the state graph has mean weight greater than 1.

\textbf{Proof of the upper bound from Proposition 3.} Let n ≥ 4 and let
S be admissible on the 3×n board. Let s\_1, \ldots, s\_k be its state
path, where k = n − 3. Summing the inequality in (b) along the path
gives

Σ over transitions (pop(new column) − 1) ≤ Φ(s\_1) − Φ(s\_k).

The total number of points is

T = {[}pop(c\_1) + \ldots{} + pop(c\_4){]} + Σ over transitions pop(new
column) = {[}pop(c\_1) + \ldots{} + pop(c\_4){]} + (n − 4) + Σ over
transitions (pop(new column) − 1) ≤ {[}pop(c\_1) + \ldots{} + pop(c\_4)
+ Φ(s\_1){]} + (n − 4) − Φ(s\_k) ≤ 6 + (n − 4) − 0 = n + 2,

using (c) for s\_1 and (a) for s\_k. ∎

\textbf{Verification of Proposition 3.} The conditions (a)--(c) were
checked by two independent implementations: one that uses the DP window
code, and one that uses only a brute-force test of all triples against
the definition of a right isosceles triangle
(\texttt{check\_certificate.py}). Both confirm (a), (b) and (c) for all
443 states and all 1736 admissible 5-tuples. The maximum of pop-sum plus
Φ over states is exactly 6, so the bound in (c) is tight for some state.

\subsection{5. Lower bound: a(3,n) ≥ n + 2}\label{lower-bound-a3n-n-2}

\textbf{Notation.} Write the block B = (0, 0, 5, 5). Mask 5 means points
in rows 0 and 2. The infinite periodic word (0,0,5,5)\^{}∞ has 2 points
in each of its 5s and none in its 0s, so each block of four columns
contains four points.

\textbf{Base sequences.} For each b ∈ \{8, 9, 10, 11\}, the following
3×b column sequence is admissible, has b+2 points, and ends with the
block (0,0,5,5):

\begin{itemize}
\tightlist
\item
  b = 8: (7, 0, 0, 7, 0, 0, 5, 5), with 10 points.
\item
  b = 9: (7, 0, 2, 1, 5, 0, 0, 5, 5), with 11 points.
\item
  b = 10: (3, 4, 1, 1, 4, 3, 0, 0, 5, 5), with 12 points.
\item
  b = 11: (6, 1, 1, 1, 1, 4, 3, 0, 0, 5, 5), with 13 points.
\end{itemize}

For n = 4, 5, 6, 7, the following sequences are admissible with n+2
points:

\begin{itemize}
\tightlist
\item
  n = 4: (7, 0, 0, 7), with 6 points.
\item
  n = 5: (7, 0, 2, 4, 5), with 7 points.
\item
  n = 6: (7, 0, 2, 2, 0, 7), with 8 points.
\item
  n = 7: (7, 0, 2, 4, 4, 4, 3), with 9 points.
\end{itemize}

\textbf{Construction.} For n ≥ 8, let b ∈ \{8, 9, 10, 11\} be the unique
value with b ≡ n (mod 4). Form the sequence consisting of base B\_b
followed by (n − b)/4 copies of the block (0,0,5,5). This has n columns
and b+2 + 4·(n−b)/4 = n+2 points.

\textbf{Admissibility.} By Lemma 2 it suffices to check the five-column
windows.

\begin{itemize}
\tightlist
\item
  Windows lying inside the base are admissible, since the base is
  admissible.
\item
  Windows lying inside the periodic part are windows of (0,0,5,5)\^{}∞.
  There are four phases, namely (0,0,5,5,0), (0,5,5,0,0), (5,5,0,0,5)
  and (5,0,0,5,5). Each is admissible.
\item
  Any window that is not of the first two kinds must start among the
  last four columns of the base. Those four columns are (0,0,5,5), which
  is the first four columns of the periodic word. So such a window is
  also a window of (0,0,5,5)\^{}∞, and it is admissible.
\end{itemize}

Hence the constructed set is admissible. It has n+2 points. ∎

The construction was also checked directly, by brute force on all
triples of points, for every n from 4 to 60 (\texttt{construct3.py},
\texttt{construct3.json}).

\subsection{6. Conclusion}\label{conclusion}

Combining Sections 4 and 5, for all n ≥ 4:

\textbf{a(3,n) = n + 2.}

For n ≤ 3 the values are a(3,1) = 3 and a(3,3) = 4 (computed).

\subsection{7. Status and limitations}\label{status-and-limitations}

\begin{itemize}
\tightlist
\item
  \textbf{Hand-checked:} Lemma 1, Lemma 2, the upper-bound deduction
  from Proposition 3, all base sequences and the periodic argument for
  the lower bound.
\item
  \textbf{Computer-assisted:} the existence and verification of the
  potential Φ in Proposition 3. This is a finite check on 443 states and
  1736 transitions. We checked it with two independent codes. A reader
  who wants a purely human proof would need a closed-form potential,
  which we did not find.
\item
  \textbf{Also computer-checked:} the lower-bound construction for n ≤
  60, and exact values up to n = 60 from the dynamic program, which
  agree with n+2 throughout.
\item
  \textbf{Literature:} the only related source we found is the OEIS
  entry for the square-grid problem, A271906 (n ≤ 11), and a 2026
  preprint on the square problem (Károlyi--Solymosi, arXiv:2607.22828).
  Neither contains this result. We did not carry out a complete
  literature search, so novelty is not guaranteed.
\end{itemize}

\subsection{8. Files}\label{files}

All files are in \texttt{./}:

\begin{itemize}
\tightlist
\item
  \texttt{dp3.py}: window checks and the dynamic program for a(3,n).
\item
  \texttt{potential3.py}: longest-path potential, D = 6.
\item
  \texttt{phi3.json}: the potential table Φ.
\item
  \texttt{check\_certificate.py}: independent verification of
  Proposition 3.
\item
  \texttt{lower3.py}, \texttt{construct3.py}, \texttt{construct3.json}:
  base sequences and the construction.
\end{itemize}

\end{document}
